\documentclass{beamer}
\usepackage{graphicx}
\usepackage{booktabs}
\usepackage{tabularx}
\usepackage{array}

\usetheme{Madrid}
\setbeamertemplate{navigation symbols}{}

\title{SDR Monitoring Spectrum System}
\subtitle{Predesign Topics Review}
\author{DesignSensorSpectrum}
\date{\today}

\begin{document}

\begin{frame}
    \titlepage
\end{frame}

\begin{frame}{Predesign Context}
    \begin{block}{System Goal}
        Build a monitoring node around an embedded compute unit, an SDR front-end, controlled RF switching, and optional remote connectivity plus geolocalization.
    \end{block}

    \vspace{0.5em}

    \begin{itemize}
        \item \textbf{SDR front-end:} Ettus USRP B200mini.
        \item \textbf{Compute unit options:} Jetson Nano or Raspberry Pi 5.
        \item \textbf{RF switching:} HMC849 SPDT for branch selection.
        \item \textbf{Remote node functions:} LTE backhaul and GNSS through SIM7600G-H.
        \item \textbf{Critical engineering axes:} power integrity, impedance control, thermals, and RF input protection.
    \end{itemize}
\end{frame}

\begin{frame}{Power Tree: How to Do It}
    \scriptsize
    \begin{enumerate}
        \item \textbf{List all loads:} compute board, SDR, RF switch, modem, fan, storage, and optional active antennas.
        \item \textbf{Separate power domains:} noisy digital loads, RF/control loads, and high-peak loads such as LTE.
        \item \textbf{Size worst-case current:} streaming SDR + modem TX + fan + storage is the design point, not idle mode.
        \item \textbf{Choose regulation:} use the official SBC input rail as the root, then derive clean 5 V / 3.3 V branches as needed.
        \item \textbf{Add protection:} fuse or current limit, reverse-polarity protection, filtering, TVS where appropriate, and ground strategy.
        \item \textbf{Validate margins:} check cable drop, peak current, boot current, and thermal rise under simultaneous load.
    \end{enumerate}

    \vspace{0.5em}

    \begin{block}{Applied to This Design}
        Root input $\rightarrow$ SBC rail $\rightarrow$ USB 3.0 SDR branch + RF switch control rail + LTE/GNSS branch + cooling branch.
    \end{block}
\end{frame}

\begin{frame}{SDR Internal Specs: Hardware and Software}
    \scriptsize
    \begin{columns}[T]
        \column{0.49\textwidth}
        \begin{block}{Hardware View}
            \begin{itemize}
                \item AD9364-based RF transceiver.
                \item \textbf{70 MHz to 6 GHz} operating range.
                \item \textbf{12-bit} ADC/DAC.
                \item Up to \textbf{61.44 MS/s}.
                \item \textbf{J1 TRX} shared TX/RX and \textbf{J2 RX2} dedicated RX.
                \item External \textbf{10 MHz / PPS} support on J3.
                \item Maximum safe RF input: \textbf{-15 dBm}.
            \end{itemize}
        \end{block}

        \column{0.49\textwidth}
        \begin{block}{Software View}
            \begin{itemize}
                \item Managed through \textbf{UHD}.
                \item Key parameters: \texttt{sample\_rate}, bandwidth, RX/TX gain, AGC, \texttt{master\_clock\_rate}.
                \item Analog LPF and digital filters are configured through UHD.
                \item IQ imbalance and DC offset correction are already built into the transceiver/UHD stack.
                \item System bottlenecks are usually host CPU, USB throughput, and buffer overrun/underrun handling.
            \end{itemize}
        \end{block}
    \end{columns}
\end{frame}

\begin{frame}{Impedance Matching: 50 $\Omega$ RF vs 75 $\Omega$ TDT}
    \small
    \begin{columns}[T]
        \column{0.48\textwidth}
        \begin{block}{50 $\Omega$ Domain}
            \begin{itemize}
                \item SDR ports, SMA cables, attenuators, switch modules, filters, and lab accessories.
                \item This should be the controlled impedance of the internal RF chain.
                \item Keep unused TX or test ports terminated with \textbf{50 $\Omega$}.
            \end{itemize}
        \end{block}

        \column{0.48\textwidth}
        \begin{block}{75 $\Omega$ TDT Domain}
            \begin{itemize}
                \item Typical for terrestrial TV antennas, splitters, and domestic coax runs.
                \item If the system connects to TDT infrastructure, adapt the transition into the SDR chain.
                \item For calibrated work, use a \textbf{$75\,\Omega \leftrightarrow 50\,\Omega$} pad, transformer, or matching section.
            \end{itemize}
        \end{block}
    \end{columns}

    \vspace{0.5em}

    \begin{alertblock}{Practical Rule}
        A direct 75 $\Omega$ to 50 $\Omega$ connection often still works, but it is a mismatch. For repeatable measurements, define one controlled adaptation point and keep the rest of the chain at 50 $\Omega$.
    \end{alertblock}
\end{frame}

\begin{frame}{Compute Unit: Expected Resources}
    \scriptsize
    \begin{tabularx}{\textwidth}{>{\raggedright\arraybackslash}p{0.24\textwidth} >{\raggedright\arraybackslash}p{0.34\textwidth} X}
        \toprule
        \textbf{Platform} & \textbf{Official Core Specs} & \textbf{Predesign Reading} \\
        \midrule
        Jetson Nano & 4x Cortex-A57 @ 1.43 GHz, 4 GB LPDDR4, 128-core Maxwell GPU, 472 GFLOPS, 5--10 W & Better when the node also needs CUDA or AI inference near the SDR pipeline. Limited CPU compared with newer SBCs. \\
        Raspberry Pi 5 & 4x Cortex-A76 @ 2.4 GHz, LPDDR4X options from 1 GB to 16 GB, VideoCore VII, 2x USB 3.0, GbE, 5 V / 5 A recommended & Stronger general CPU platform for host-side DSP, logging, and orchestration if GPU acceleration is not the main requirement. \\
        \bottomrule
    \end{tabularx}

    \vspace{0.5em}

    \begin{block}{Sizing Guideline}
        For plain spectrum monitoring and control, Raspberry Pi 5 is the stronger CPU-first option. For AI-assisted classification, camera fusion, or CUDA workflows, Jetson Nano keeps an advantage at the platform level.
    \end{block}
\end{frame}

\begin{frame}{Thermal Throttling}
    \small
    \begin{itemize}
        \item Thermal throttling is not only a CPU problem; it impacts \textbf{sample stability}, USB reliability, modem uptime, and enclosure behavior.
        \item Main heat sources in this design are:
        \begin{itemize}
            \item compute SoC,
            \item LTE modem during transmission,
            \item storage and power stages,
            \item confined airflow around the SDR host.
        \end{itemize}
        \item Raspberry Pi 5 official operating temperature is \textbf{0\textdegree C to 70\textdegree C}; the enclosure should be designed to stay below thermal reduction zones.
        \item Jetson Nano power modes are low, but sustained AI or video tasks still need heatsinking and airflow.
        \item Predesign mitigation: fan header budget, vent path, modem spacing, thermal pads/heatsinks, and logging of temperature under real SDR streaming load.
    \end{itemize}
\end{frame}

\begin{frame}{Undervoltage}
    \small
    \begin{itemize}
        \item Undervoltage causes unstable USB peripherals, SDR disconnects, sample drops, storage errors, and LTE modem resets.
        \item Worst cases appear during \textbf{boot}, \textbf{USB 3.0 traffic}, and \textbf{SIM7600 TX current peaks}.
        \item Raspberry Pi 5 official guidance:
        \begin{itemize}
            \item it can boot from \textbf{5 V / 3 A},
            \item but high-power peripheral current is then restricted,
            \item \textbf{5 V / 5 A} gives more USB current budget and more system margin.
        \end{itemize}
        \item Predesign rules:
        \begin{itemize}
            \item keep input cables short and low resistance,
            \item avoid powering modem peaks from a weak shared branch,
            \item leave current margin for fan and storage,
            \item validate supply sag with all subsystems active at once.
        \end{itemize}
    \end{itemize}
\end{frame}

\begin{frame}{Internet or Node Connection}
    \small
    \begin{columns}[T]
        \column{0.48\textwidth}
        \begin{block}{Lab / Fixed Node}
            \begin{itemize}
                \item Gigabit Ethernet for stable remote access.
                \item SSH, rsync, VPN, or web dashboards.
                \item Easier for continuous logging and software updates.
            \end{itemize}
        \end{block}

        \column{0.48\textwidth}
        \begin{block}{Field / Mobile Node}
            \begin{itemize}
                \item SIM7600G-H provides LTE backhaul.
                \item USB is the recommended integration path with Jetson-class hosts.
                \item Good for remote campaigns, telemetry, and alerting.
            \end{itemize}
        \end{block}
    \end{columns}

    \vspace{0.5em}

    \begin{alertblock}{Node Architecture Guideline}
        Do local capture and local reduction first. Send metadata, alarms, and summarized spectra over the network; send raw wideband IQ only when the link budget allows it.
    \end{alertblock}
\end{frame}

\begin{frame}{RF Switching}
    \small
    \begin{itemize}
        \item The \textbf{HMC849 SPDT} is the predesign choice for switching between two RF branches.
        \item Current architecture idea:
        \begin{itemize}
            \item USRP \textbf{RX2} connected to the common switch port,
            \item output A and output B connected to two antennas or two filter paths,
            \item GPIO from the compute unit selects the active branch.
        \end{itemize}
        \item Typical uses:
        \begin{itemize}
            \item antenna diversity,
            \item band-select filtering,
            \item protected bypass topologies.
        \end{itemize}
        \item The switch is \textbf{not} an attenuator, limiter, or surge protector.
        \item Predesign checks: logic threshold compatibility, insertion loss, isolation, RF power limit, and board-level grounding.
    \end{itemize}
\end{frame}

\begin{frame}{Geolocalization}
    \small
    \begin{itemize}
        \item The \textbf{SIM7600G-H} adds \textbf{GPS, BeiDou, GLONASS, Galileo, and QZSS} support.
        \item Geolocation is useful for:
        \begin{itemize}
            \item tagging spectrum captures by position,
            \item mobile drive-test style campaigns,
            \item asset tracking of deployed nodes,
            \item correlating events with location and route.
        \end{itemize}
        \item Integration options:
        \begin{itemize}
            \item GNSS queried over AT commands,
            \item stored in local logs with timestamp and node ID,
            \item forwarded to a server together with alarms or summary measurements.
        \end{itemize}
        \item Important distinction: GNSS gives position; precise SDR coherence still benefits from a dedicated \textbf{10 MHz / PPS} timing strategy if future synchronization is required.
    \end{itemize}
\end{frame}

\begin{frame}{Predesign Output}
    \begin{itemize}
        \item Keep the internal RF chain at \textbf{50 $\Omega$} and define explicit adaptation for TDT-style 75 $\Omega$ inputs.
        \item Size power around \textbf{peak simultaneous load}, not around nominal board values.
        \item Pick \textbf{Raspberry Pi 5} for CPU-centric SDR nodes, or \textbf{Jetson Nano} when GPU/AI integration is part of the roadmap.
        \item Treat thermal throttling and undervoltage as first-order system risks.
        \item Use LTE and GNSS as node services, not as substitutes for RF timing or front-end protection.
    \end{itemize}
\end{frame}

\end{document}
